\documentclass[../main.tex]{subfiles}

\begin{document}

\section{Testes com Exercicios do Livro}

Os testes seguem os exercicios indicados do livro de Neide Franco. As entradas
ficam em \texttt{input/exercicios} e as saidas em CSV ficam em
\texttt{output/exercicios}.

\subsection{Exercicios do Capitulo 8}

No exercicio 8.1, foi usada regressao linear para o numero de acidentes e
regressao quadratica para os acidentes por 10000 veiculos. Seguindo a
observacao do enunciado, o tempo foi deslocado usando \(x = ano - 1980\).

No exercicio 8.5, os dados da placa de orificio foram ajustados por um
polinomio de grau 2. No exercicio 8.11, o modelo \(a x^b\) foi linearizado com
logaritmos para ajustar separadamente os PNBs em dolares correntes e constantes.

\subsection{Exercicios do Capitulo 10}

No exercicio 10.2, os tres pontos da trajetoria do projetil geram o polinomio
\[
    P(x) = x - 0.04x^2.
\]
A partir dele foram calculados o angulo de lancamento, a velocidade inicial e a
altura em \(x=5\).

Nos exercicios 10.6 e 10.9 foram usadas interpolacoes polinomiais. O exercicio
10.6 calcula resistencias por polinomios de grau 2 e 3 usando os pontos mais
proximos de cada consulta. O exercicio 10.9 usa interpolacao direta para
calcular a corrente no varistor e interpolacao inversa de grau 2 para estimar a
tensao quando \(E_1=10\).

\subsection{Exercicios do Capitulo 11}

Os exercicios 11.1 e 11.6 foram resolvidos por Simpson composto com refinamento
automatico ate a variacao relativa ficar abaixo de \(10^{-5}\). O exercicio
11.11 calcula a area acumulada de uma distribuicao normal ate \(x=5\), tambem
por integracao numerica.

\end{document}
